% Source: https://artofproblemsolving.com/community/c3h1503732_new_problem
% Stable identity: aops-topic-1503732-q03
Number the positions $1$ to $9$ around the circle in the given order, and let $x_j$ be the number at position $j$. Consider
\[
T=1\cdot x_1+2\cdot x_2+\cdots+9\cdot x_9 \pmod 9 .
\]
A move acts on three consecutive positions. Their weights are congruent modulo $9$ to $j$, $j+1$, $j+2$ for some $j$ (also when the three positions wrap around from $9$ to $1$, since $9\equiv0$). Both orientations of the triple give the same change modulo $9$:
\[
j-2(j+1)+(j+2)=0\qquad\text{or}\qquad (j+2)-2(j+1)+j=0 .
\]
So $T$ modulo $9$ is invariant.

At the start,
\[
T=1+8+15+32+30+12+21+72+63=254\equiv2\pmod9 .
\]
If all numbers were $5$, we would have $T=5(1+2+\cdots+9)=225\equiv0\pmod9$. Since $2\not\equiv0$, this can never happen: the answer is $\boxed{\text{no}}$.
